\section{Generic symmetric monodromy and closure genus}\label{sec:generic}
\begin{theorem}\label{thm:generic}
For every \(n\ge2\), the family \eqref{eq:family} contains a nonempty Zariski-open subset of degree-\(n\) maps with \(2n-2\) simple critical points and distinct critical values. On this subset,
\begin{equation}\label{eq:main}
 \boxed{\Mon(F_n)=S_n,\qquad
 g(\widetilde F_n)=1+\frac{n!}{2}(n-3).}
\end{equation}
For \(n\ge3\), the deck group is trivial and the residual equal-value correspondence has degree \(n-1\) over \(\CC(y)\) and is irreducible. These statements concern the unchanged target \(v=F_n(y)\).
\end{theorem}

\subsection{A construction proving the generic locus is nonempty}
It is not enough to invoke generic rational maps: coefficient reversal imposes a constraint, and one must show that this constraint does not force repeated branch values. The following construction does so inside the odd coordinate chart.

Start with \(H_1(z)=z\) for odd degrees and \(H_2(z)=z/(z^2+1)\) for even degrees. The first has no critical points; the second has simple critical points \(\pm1\), with values \(\pm1/2\). Both are unramified at zero and infinity. Their finite poles, if any, are simple, and their asymptotic behavior is respectively \(cz\) and \(c/z\), with \(c\ne0\).

Suppose \(H_d\) has these properties, with all critical points simple and all critical values distinct. Choose a finite \(b\ne0\) such that \(\pm b\) are neither poles nor critical points and \(H_d(b)\) is finite, nonzero, and different from every old critical value and its negative. Only finitely many choices are forbidden. For small \(\varepsilon\ne0\), define
\begin{equation}\label{eq:insertion}
 H_{d+2}(z)=H_d(z)+\varepsilon\frac{z}{z^2-b^2}.
\end{equation}
It remains odd, adds simple poles at \(\pm b\), and has degree \(d+2\). The old poles keep their nonzero residues. The numerator cannot cancel a new pole because its residue is \(\varepsilon/2\). The behavior at infinity retains its type and its nonzero leading coefficient for sufficiently small \(\varepsilon\).

Write \(a=H_d'(b)\ne0\) and \(z=b+\xi\). The new pole contributes \(\varepsilon/(2\xi)+O(\varepsilon)\). Solving the derivative equation by a local square-root expansion gives two new critical points and their values:
\begin{align}\label{eq:local_insertion}
 z&=b\pm\sqrt{\varepsilon/(2a)}+O(\varepsilon),\notag\\
 H_{d+2}(z)&=H_d(b)\pm\sqrt{2\varepsilon a}+O(\varepsilon).
\end{align}
For a rigorous local parameter, put \(\varepsilon=t_0^2\) and \(\xi=t_0\zeta\); after multiplication by \(\zeta^2\), the limiting derivative equation has two simple roots \(\zeta=\pm(2a)^{-1/2}\). The implicit-function theorem gives the displayed branches. Their critical points are simple. Oddness supplies the two opposite branches near \(-b\). Since \(H_d(b)\ne0\), these four values are distinct for sufficiently small nonzero \(t_0\) and avoid the old values.

Old simple critical points persist by the same theorem. Nonzero local linear coefficients at zero and infinity persist. We have exhibited
\((2d-2)+4=2(d+2)-2\) simple critical points. Riemann--Hurwitz leaves no additional ramification. Induction therefore constructs the required odd maps in every degree.

To return to the precise coefficient chart, conjugate by \(\mathcal C^{-1}\). A generic nonzero target scaling of \(H\) avoids the condition that the resulting numerator vanish at \(y=0\); it preserves oddness and branch simplicity. Reciprocity makes the resulting homogeneous numerator and reversed denominator proportional. At \(y=1\), the value is \(+1\), with no cancellation, so the proportionality sign is \(+1\). Scaling the pair to make the numerator constant term one gives \(Q=P^*\) exactly. The parity types in \eqref{eq:odd_types} give the required degree and value at \(-1\). Thus these are witnesses in \eqref{eq:family}, not in an unrelated larger family.

Coprimeness, simple critical divisor, and distinct critical values are open algebraic conditions: their failures are detected by homogeneous resultants and discriminants. The coefficient space \(\CC^n\) is irreducible. Existence of a witness in each degree consequently proves that the desired locus is nonempty and Zariski open.

\subsection{Monodromy, partners, and genus}
At a simple branch value precisely one sheet pair is exchanged. All branch cycles are therefore transpositions, and they generate the monodromy group. Connectedness of the source makes that action transitive. A transitive group generated by transpositions is \(S_n\): form the graph whose edges are those transpositions; it is connected, and transpositions along a connected graph generate every permutation.

Deck permutations commute with monodromy on a regular fiber. The centralizer of the natural \(S_n\) action is trivial for \(n\ge3\). After choosing one sheet, its stabilizer is \(S_{n-1}\), which acts transitively on the other \(n-1\) sheets. The residual equal-value correspondence is therefore irreducible over the chosen source field. A rational partner \(M\) would satisfy \(F_n\circ M=F_n\); multiplicativity of degree forces \(\deg M=1\), making it a Möbius deck transformation. None exists generically for \(n\ge3\).

Finally the Galois closure has degree \(n!\). It has \(2n-2\) branch values, each with inertia order two. Galois Riemann--Hurwitz yields
\begin{equation}\label{eq:RH_general}
 2g-2=n!\left[-2+(2n-2)\left(1-\frac12\right)\right]
       =n!(n-3).
\end{equation}
This proves Theorem~\ref{thm:generic}. The construction proves all degrees; no extrapolation from the two numerical witnesses is used.

\subsection{Why the quadratic member is exceptional}\label{sec:quadratic}
A separable degree-two extension is automatically Galois. Geometrically there is exactly one other point in a generic fiber, so the partner cannot be permuted with competing choices. It is a rational automorphism of the sphere and hence Möbius. For the uncanceled quadratic Three-Way map,
\begin{equation}\label{eq:quadratic_partner}
 T(y)=\frac{(A+1)y-B}{By-(A+1)},\qquad F_2(T(y))=F_2(y).
\end{equation}
The degree-two condition is \(A\ne1\) and \((A+1)^2-B^2\ne0\). The latter is also the nonzero determinant condition for \(T\). Reciprocity \(R(y)=1/y\) normalizes the unique nonidentity deck element, so it commutes with \(T\). They are distinct and generate \(V_4\). Its subgroup \(\{1,T\}\) fixes \(F_2\); its other two elements invert \(F_2\). Thus \(V_4\) fixes the further target quotient, for example \(F_2+1/F_2\), while \(\Deck(F_2)=C_2\).

The following table always concerns \(F_n\) itself, after cancellation and on its generic locus. The lock column assumes no cancellation at the indicated points.
\begin{table}[htbp]\centering\small
\begin{tabularx}{\linewidth}{L{.25\linewidth}YYY}\toprule
Property & \(n=2\) & \(n=3\) & \(n=4\)\\\midrule
Degree & 2 & 3 & 4\\
Reciprocity & \(F(1/y)=1/F(y)\) & same & same\\
Other-sheet equation & linear, rational partner & irreducible quadratic & irreducible cubic\\
Deck group & \(C_2\) & trivial & trivial\\
Monodromy & \(C_2=S_2\) & \(S_3\) & \(S_4\)\\
Closure degree; genus & 2; 0 & 6; 1 & 24; 13\\
Forced \(+1\) locks & \(y=\pm1\) & \(y=1\) & \(y=\pm1\)\\
Selected exceptions & cancellation / constant & cyclic, nodal, cancellation & decomposable, cyclic, \(V_4\), cancellation\\\bottomrule
\end{tabularx}
\caption{Generic hierarchy over the unchanged rational-map target. Even parity restores a lock at \(-1\), but does not restore the quadratic partner mechanism.}\label{tab:hierarchy}
\end{table}
\fig{01_hierarchy}{Degrees and closure genera refer to three distinct extensions over their own unchanged targets. Paper I's elliptic curve is associated with the separate larger quadratic tower. The lower row records the residual other-sheet degree.}{fig:hierarchy}
